\section{Build an encoder from a correct literal baseline}
Start with a representation that writes one RGBA chunk for every input pixel. It
requires no predictor or cache lookup and gives a simple upper bound: five payload
bytes per pixel. This baseline is useful for checking the header, component order,
marker and exact written length before introducing compression decisions.

Next keep the previous pixel and accumulate repeated pixels into runs. When a different
pixel arrives, flush the pending run before describing the new pixel. If the cache
contains the entire new RGBA value, an INDEX chunk can supply it. Otherwise compare
alpha with the previous alpha. Changed alpha requires RGBA. Unchanged alpha permits
DIFF when all its small changes fit, then LUMA when its change and residuals fit, and
otherwise RGB. This order is a concrete encoder policy to study, not the only legal
way to represent an image.

Update the encoder cache to agree with the decoder state produced by the bytes already
emitted. A hash collision is not a match. A run must be emitted at length sixty-two
or the final pixel, and a pending shorter run must precede a following different color.
An output count that ignores a flushed run points the next write at the wrong position.
Keep the pending run length and byte cursor separate.

The completed exercise accepts different valid encodings that recover the exact input
pixels and preserve the supplied descriptive flags. It also checks compression on
uniform images. Otherwise an encoder that always emitted RGBA would pass a lossless
round trip without meeting the run-encoding objective. For initial opaque black, use
no more than \(22+\lceil n/62\rceil\) bytes. For another uniform pixel, allow
\(27+\lceil(n-1)/62\rceil\): one at-most-five-byte first pixel followed by runs.
No exact instruction count or whole-file byte identity is a grading requirement.
